\chapter{Contexte et état de l'art}
\label{chap:contexte}

\noindent\itshape
Ce chapitre situe le problème — la sécurité des applications fondées sur les LLM et son évaluation — puis dresse l'état de l'art des vulnérabilités, des méthodes d'attaque automatisées et des outils, pour faire apparaître le verrou méthodologique que ce mémoire entend lever.
\normalfont
\vspace{0.5cm}

\section{Contexte général}

Les grands modèles de langage (LLM) animent aujourd'hui des produits utilisés au quotidien~: agents conversationnels, génération augmentée par récupération (RAG), et surtout \textbf{agents} qui appellent des outils, interrogent des bases ou déclenchent des actions. Cette intégration a fait surgir une surface d'attaque d'un type nouveau. L'entrée du système est du \textbf{langage naturel}~; le comportement est \textbf{non déterministe}~; et, fait décisif, les instructions de confiance (le prompt système) partagent la même fenêtre de contexte que des données fournies par l'utilisateur, parfois hostiles.

La sécurité applicative classique repose sur des frontières nettes — entre code et données, entre requête et réponse. Une application LLM \emph{n'a pas} cette séparation forte~: rien, au niveau du modèle, ne distingue une consigne légitime d'une consigne injectée dans un document. C'est la racine d'une famille de vulnérabilités propres aux LLM, que les contrôles traditionnels (pare-feu, validation de schéma) ne couvrent pas.

\section{La difficulté d'évaluer la sécurité d'une application LLM}

Tester une application LLM se heurte à trois obstacles.

D'abord le \textbf{non-déterminisme}~: une même attaque peut réussir un essai sur trois. Le succès n'est pas binaire mais \emph{probabiliste}, ce qui impose de raisonner en taux et en intervalles de confiance plutôt qu'en simple « ça marche / ça ne marche pas ». Ensuite, l'\textbf{espace d'attaque en langage naturel} est immense~: là où une injection SQL se décline en un nombre fini de motifs, une attaque par prompt se reformule, se traduit, s'encode à l'infini. Enfin, la \textbf{définition du succès} est souvent affaire de jugement~: « cette réponse constitue-t-elle une fuite~?» n'a pas toujours de réponse mécanique. D'où la tentation de confier l'arbitrage à un autre LLM — dont la fiabilité, justement, est rarement vérifiée.

Face à l'ampleur de cet espace, le red teaming des applications LLM doit être \textbf{automatisé}, au moins en partie. C'est l'objet des méthodes et des outils décrits plus loin.

\section{Le problème adressé par ce travail}

Deux manques motivent ce mémoire.

Le premier est un manque d'\textbf{accessibilité}. Les frameworks de red teaming existants sont puissants, mais lourds~: courbe d'apprentissage, configuration experte, dépendances nombreuses. Il manque un outil léger, mono-conteneur, fondé sur le principe \emph{Bring Your Own Key} (l'utilisateur fournit ses clés, conservées en mémoire seulement), observable, et compréhensible par un public non spécialiste, qui rattache ses tests à un référentiel reconnu.

Le second manque, plus profond, concerne la \textbf{rigueur de l'évaluation des outils eux-mêmes}. C'est le verrou central de ce travail. La première version du projet illustrait deux faiblesses méthodologiques courantes dans le domaine~:
\begin{itemize}
  \item une \textbf{validation circulaire}~: l'outil n'était éprouvé que sur des cibles que l'auteur avait lui-même rendues vulnérables, sans aucune référence indépendante~;
  \item un \textbf{juge LLM jamais validé}~: tous les résultats reposaient sur le verdict d'un modèle dont la justesse n'avait jamais été mesurée. Le juge était l'oracle implicite, sans contrôle.
\end{itemize}

\section{Taxonomie des vulnérabilités : l'OWASP Top~10 for LLM Applications (2025)}

Le référentiel de fait pour la sécurité des applications LLM est l'\textbf{OWASP Top~10 for LLM Applications}, dont l'édition 2025 compte dix catégories (tableau~\ref{tab:owasp}).

\begin{table}[H]
\centering
\caption{OWASP Top~10 for LLM Applications (édition 2025)}
\label{tab:owasp}
\begin{tabularx}{\textwidth}{l X}
\toprule
\textbf{Code} & \textbf{Catégorie} \\
\midrule
LLM01 & Prompt Injection (injection de prompt) \\
LLM02 & Sensitive Information Disclosure (divulgation d'information sensible) \\
LLM03 & Supply Chain (chaîne d'approvisionnement) \\
LLM04 & Data and Model Poisoning (empoisonnement de données / modèle) \\
LLM05 & Improper Output Handling (mauvaise gestion des sorties) \\
LLM06 & Excessive Agency (agentivité excessive) \\
LLM07 & System Prompt Leakage (fuite de prompt système) \\
LLM08 & Vector and Embedding Weaknesses (faiblesses des vecteurs / \emph{embeddings}) \\
LLM09 & Misinformation (désinformation) \\
LLM10 & Unbounded Consumption (consommation non maîtrisée) \\
\bottomrule
\end{tabularx}
\end{table}

L'édition 2025 introduit notamment \emph{System Prompt Leakage} (LLM07), \emph{Vector and Embedding Weaknesses} (LLM08) et \emph{Unbounded Consumption} (LLM10), et absorbe l'ancienne \emph{Overreliance} dans \emph{Misinformation} (LLM09).

RedAgent cible les catégories \textbf{testables en boîte noire, via une API JSON, et au succès objectivement décidable}~: \textbf{LLM01}, \textbf{LLM02}, \textbf{LLM06} et \textbf{LLM07}. Les autres sont documentées comme hors de ce cadre~: LLM03/04 relèvent de la phase d'entraînement~; LLM05 se joue côté intégration applicative~; LLM08 est interne au RAG~; LLM09 suppose un jugement de véracité sans flag objectif~; LLM10 relève du déni de service. Ce périmètre n'est pas une commodité~: il garantit que chaque succès mesuré l'est contre une \emph{vérité-terrain} sans ambiguïté.

\section{Méthodes d'attaque automatisées}

On distingue deux grandes familles.

\paragraph{Attaques à accès interne (\emph{white-box}).} L'approche \textbf{GCG} (\emph{Greedy Coordinate Gradient}, Zou et al., 2023) optimise par gradient un suffixe adversarial qui maximise la probabilité d'une réponse nuisible. Redoutable, mais elle exige l'accès aux poids du modèle — inapplicable à une cible de production en boîte noire.

\paragraph{Attaques en boîte noire pilotées par LLM.} Elles emploient un \emph{LLM attaquant} qui raffine ses prompts à partir des réponses de la cible, sans gradient~: \textbf{PAIR} (\emph{Prompt Automatic Iterative Refinement}, Chao et al., 2023) introduit la boucle attaquant–cible avec retour~; \textbf{TAP} (\emph{Tree of Attacks with Pruning}, Mehrotra et al., 2024) l'étend à une recherche arborescente avec élagage~; \textbf{Crescendo} (Russinovich et al., 2024) escalade progressivement, sur plusieurs tours, d'une demande bénigne vers l'objectif.

\paragraph{Positionnement de RedAgent.} RedAgent opère en boîte noire et combine des \textbf{graines génériques} (une par classe OWASP~: vidage de configuration, « répète ce qui précède », TLDR, cadrage d'autorité) avec des \textbf{mutations générées par un LLM attaquant}. C'est une approche \emph{semi-autonome}, délibérément plus simple que TAP ou PAIR. Ce choix est assumé~: l'objet du mémoire n'est pas la sophistication de l'attaque, mais la rigueur avec laquelle on en mesure le succès.

\section{Outils et frameworks de red teaming}

Plusieurs cadres outillés existent (tableau~\ref{tab:outils}). Ils sont matures, mais conçus pour l'ingénieur sécurité, avec l'installation et les abstractions que cela suppose.

\begin{table}[H]
\centering
\caption{Principaux outils de red teaming LLM}
\label{tab:outils}
\begin{tabularx}{\textwidth}{l l X}
\toprule
\textbf{Outil} & \textbf{Éditeur} & \textbf{Profil} \\
\midrule
PyRIT & Microsoft & Framework extensible~: orchestrateurs multi-tours, convertisseurs, scorers, mémoire. \\
garak & NVIDIA & Scanner à \emph{probes} (familles d'attaques prédéfinies). \\
promptfoo & open source & Tests et évaluations de prompts, orienté CI/CD. \\
Giskard / GOAT & Giskard & Red teaming automatisé multi-tours. \\
\bottomrule
\end{tabularx}
\end{table}

Le créneau de RedAgent est complémentaire et non concurrent~: léger, mono-conteneur, BYOK, observable en direct, aligné OWASP, pensé pour un public non expert. Il ne prétend pas remplacer PyRIT ou garak, mais offrir un point d'entrée accessible — et porter une contribution de \emph{méthode}.

\section{Bancs d'évaluation et verrou de la mesure}

Évaluer un outil suppose une vérité-terrain du succès. On dispose pour cela de ressources comme \textbf{Gandalf} (Lakera), jeu public d'injection de prompt où l'extraction d'un mot de passe est confirmée par un endpoint dédié~; de \textbf{CTF OWASP-LLM} à flags déterministes~; ou de jeux de prompts nuisibles tels que HarmBench ou AdvBench — dont le succès, hélas, est souvent évalué par un classifieur ou un juge LLM.

\paragraph{Le juge LLM : une fiabilité étudiée ailleurs, pas ici.} Il serait inexact de prétendre que la fiabilité des juges LLM n'a jamais été étudiée. Un corpus déjà substantiel l'examine dans le cadre du \emph{LLM-as-a-judge} appliqué à l'évaluation de qualité~: Zheng et al. (2023), avec MT-Bench et Chatbot Arena, mesurent l'accord entre un juge LLM et des préférences humaines (de l'ordre de 80~\%, comparable à l'accord inter-humains) tout en documentant des biais systématiques — biais de position, de verbosité, d'auto-complaisance~; des travaux de synthèse ultérieurs (Gu et al., 2024) recensent ces effets et les protocoles d'atténuation. \textbf{Mais ce corpus mesure le juge sur une tâche de \emph{préférence} ou de \emph{notation de qualité}} (« quelle réponse est la meilleure~?», « cette réponse est-elle utile~?»), dans un cadre de \emph{benchmark}. Il ne porte presque jamais sur l'usage qui nous occupe~: l'\textbf{arbitrage binaire du succès d'une attaque} en red teaming d'\emph{outils} (« cette réponse constitue-t-elle, oui ou non, une vulnérabilité~?»), et plus rarement encore contre une vérité-terrain humaine \emph{établie pour cet arbitrage précis}. Le verrou n'est donc pas que la question soit vierge, mais qu'elle reste ouverte là où on en a besoin~: au point exact où un juge LLM décide, en pratique, si un outil de sécurité a « réussi ».

C'est ici qu'apparaît le \textbf{verrou méthodologique}. La littérature multiplie les attaques et les outils, mais traite la \emph{mesure} du succès — quand c'est un juge LLM qui la rend — comme acquise dans ce contexte précis. Les deux écueils déjà cités — validation circulaire et juge non validé pour l'arbitrage de succès — y sont fréquents. La contribution de ce mémoire se situe précisément à cet endroit~: rendre le succès \emph{objectivement} décidable par un oracle déterministe, et rendre le juge LLM \emph{mesurable} contre cet oracle puis contre un jugement humain. Le chapitre suivant en présente la conception.
